\documentclass[12pt,a4paper]{article}

% ============================================================
%  FA-Theorem — Federated Audit Theorem:
%  Zero-Knowledge Multi-Party Audit Completeness via
%  Situs Homomorphic Composition
%  arXiv-ready · English · Standalone (SCX-citable, not dependent)
% ============================================================

\usepackage[utf8]{inputenc}
\usepackage[T1]{fontenc}
\usepackage{amsmath,amssymb,amsfonts}
\usepackage{mathtools}
\usepackage{amsthm}
\usepackage{bm}
\usepackage{graphicx}
\usepackage{xcolor}
\usepackage{hyperref}
\usepackage{booktabs}
\usepackage{enumitem}
\usepackage[numbers]{natbib}
\usepackage{algorithm}
\usepackage{algpseudocode}

% ---------- Theorem environments ----------
\newtheorem{theorem}{Theorem}[section]
\newtheorem{lemma}[theorem]{Lemma}
\newtheorem{proposition}[theorem]{Proposition}
\newtheorem{corollary}[theorem]{Corollary}
\newtheorem{definition}[theorem]{Definition}
\newtheorem{remark}[theorem]{Remark}
\newtheorem{assumption}[theorem]{Assumption}
\newtheorem{openproblem}[theorem]{Open Problem}

% ---------- Status tags ----------
\newcommand{\rigorous}{\textcolor{green!50!black}{\small\textsf{[Rigorous]}}}
\newcommand{\conditionallyrigorous}{\textcolor{orange}{\small\textsf{[Conditionally Rigorous]}}}
\newcommand{\openquest}{\textcolor{red!70!black}{\small\textsf{[Open]}}}

% ---------- Macros ----------
\newcommand{\R}{\mathbb{R}}
\newcommand{\E}{\mathbb{E}}
\newcommand{\V}{\mathbb{V}}
\newcommand{\I}{\mathbb{I}}
\newcommand{\cX}{\mathcal{X}}
\newcommand{\cY}{\mathcal{Y}}
\newcommand{\cD}{\mathcal{D}}
\newcommand{\cS}{\mathcal{S}}
\newcommand{\cP}{\mathcal{P}}
\newcommand{\cA}{\mathcal{A}}
\newcommand{\ind}[1]{\mathbf{1}_{[#1]}}
\newcommand{\argmax}{\operatorname*{arg\,max}}
\newcommand{\argmin}{\operatorname*{arg\,min}}
\newcommand{\KL}{D_{\mathrm{KL}}}
\newcommand{\Situs}{\mathrm{Situs}}
\newcommand{\Cercis}{\mathrm{Cercis}}
\newcommand{\Commit}{\mathrm{Commit}}
\newcommand{\Verify}{\mathrm{Verify}}
\newcommand{\ZK}{\mathrm{ZK}}

\title{\bfseries Federated Audit Completeness:\\
Zero-Knowledge Multi-Party Auditing\\
via Situs Homomorphic Composition}

\author{Anonymous Author(s)}

\date{\today}

\begin{document}
\maketitle

\begin{abstract}
The SCX Audit Sword establishes that any use of the Cercis Score can be
independently audited — provided the auditor possesses complete data access.
In federated scenarios where $K$ parties hold private datasets and interact
only through aggregated statistics, the Audit Sword encounters a structural
obstacle: how does one audit what one cannot see?  We formalize this as the
\emph{Federated Audit Problem} and prove three results.  (i)~When the Situs
encoding admits a homomorphic composition operator $\oplus$ such that
$\Situs(\cup_k \cD_k) \cong \oplus_k \Situs(\cD_k)$, and each party's Cercis
Score submission is verified by a zero-knowledge proof $\pi_{\ZK}$, the
federated audit achieves $\varepsilon$-equivalence to the full-data audit,
with $\varepsilon \to 0$ under vanishing differential privacy budget and
growing sample size.  The Situs homomorphic composition construction is
\textbf{conditionally rigorous} — given the operator, the $\varepsilon$-bound
is strict; the operator's existence is an \textbf{open problem}.
(ii)~Without $\ZK$ verification, the information loss is lower-bounded by
$\mathrm{const} \cdot \eta_k \cdot |\cD_k|/|\cD_{\mathrm{total}}|$, a
\textbf{rigorous} consequence of Theorem~3 and the data processing inequality.
(iii)~We outline the ``Telegraph Public Good'' extension where decentralized
storage-layer commitments replace $\ZK$ proofs with Proofs of Storage,
further tightening the $\varepsilon$ bound — an \textbf{open problem} with
a clear theoretical path.  Collectively, the results establish that federated
auditing is possible without raw data access, but only when the information
architecture provides verifiable audit fingerprints.
\end{abstract}

% ============================================================
\section{Introduction}
% ============================================================

The SCX framework~\citep{scx2024cercis} established the \textbf{Audit Sword}
principle: every use of the Cercis Score must be independently verifiable by
any third party.  This is not an optional feature — it is a structural
requirement that follows from Theorem~3 (the Noise--Difficulty
Unidentifiability Theorem).  If you cannot audit the detector, you cannot
distinguish whether it is flagging genuine label noise or merely reclassifying
difficult-but-correct samples according to its own inductive biases.

The Audit Sword, however, carries a hidden premise: that the auditor has
\textbf{complete access to all data}.  In an increasingly federated world —
hospitals that cannot share patient records across jurisdictions, financial
institutions barred from exposing transaction logs, edge devices with
bandwidth and privacy constraints — this premise fails.  The very parties
whose models most urgently require auditing (because they operate on
sensitive, high-stakes data) are precisely those who cannot grant the access
that auditing presupposes.

This creates the \textbf{Federated Audit Paradox}.  Federated learning has
developed an extensive privacy-preserving toolbox — Secure
Aggregation~\citep{bonawitz2017practical}, Differential
Privacy~\citep{abadi2016deep}, Multi-Party
Computation~\citep{evans2018pragmatic} — but each of these tools is designed
to \emph{prevent information leakage}.  The SCX requirement is the inverse:
we need to \emph{enable auditability} without compromising privacy.  The
existing toolbox solves the problem of locking the door; we need to solve
the problem of auditing through it.

\medskip\noindent
\textbf{The Core Tension.}
At the heart of the paradox lies an information asymmetry.  The auditor's
full-information statistic $T_{\mathrm{full}}$ operates on raw $(x, y)$ pairs
from the union of all parties' datasets.  The federated auditor instead
receives only \emph{audit fingerprints} — the Cercis Scores $S_k$ computed
locally by each party, potentially encrypted, along with the aggregated model
$M_\Pi$.  The question is not whether the fingerprints leak the raw data (they
should not), but whether they preserve enough statistical structure to support
audit conclusions with the same power as the full-information benchmark.

\medskip\noindent
\textbf{Contributions.}
\begin{enumerate}[label=(\arabic*)]
    \item \textbf{Federated Audit $\varepsilon$-Equivalence}
          (Theorem~\ref{thm:fed-equiv}): Under Situs homomorphic composition
          and $\ZK$ verification, $|\mathrm{Power}(T_{\mathrm{fed}}) -
          \mathrm{Power}(T_{\mathrm{full}})| \leq \varepsilon(K, \delta)$,
          with $\varepsilon \to 0$ as $n \to \infty$ and $\delta \to 0$.
          \conditionallyrigorous~/ \openquest~(Situs homomorphism)
    \item \textbf{Information Lower Bound}
          (Theorem~\ref{thm:info-lower}): Without $\ZK$ verification, audit
          power loss is bounded below by a fraction proportional to the
          unverified noise share.  \rigorous
    \item \textbf{Telegraph Public Good Extension}
          (Section~\ref{sec:telegraph}): Decentralized storage commitments
          as a $\ZK$ replacement.  \openquest
\end{enumerate}

% ============================================================
\section{Preliminaries}
% ============================================================

\subsection{SCX Audit Architecture}

We recall the essential SCX components required for federated audit analysis.
For a dataset $\cD = \{(x_i, y_i)\}_{i=1}^n$, the \textbf{Cercis Score}
$S: \cX \to \R_{\geq 0}$ decomposes as
\begin{equation}\label{eq:cercis}
    S(x) = Q(x) + \eta N(x),
\end{equation}
where $Q(x) \in [0,1]$ measures label quality (model confidence, expert
agreement) and $N(x) \in [0,1]$ measures novelty (dissimilarity to previously
audited samples).  The parameter $\eta \geq 0$ is the Sprint noise--difficulty
coupling coefficient.

The \textbf{Situs operator} encodes a data-generating distribution into a
metric-measure space:
\begin{equation}\label{eq:situs}
    \Situs(P) = (\cX, d_P, \mu_P),
\end{equation}
where $d_P$ is a task-adapted metric and $\mu_P$ is the marginal measure on
$\cX$.  Two distributions are compared via the 1-Wasserstein distance
$d_{\Situs}(P, Q) = W_1(\Situs(P), \Situs(Q))$.

\textbf{Theorem~3} (Noise--Difficulty Unidentifiability).  From Cercis Score
observations alone, no algorithm can distinguish label noise from intrinsic
difficulty with power exceeding its significance level.  Formally, for any
test $\psi$ based on $\{S(x_i)\}_{i=1}^n$,
$\sup_\psi |\E_{H_1}[\psi] - \E_{H_0}[\psi]| \leq \alpha + o_n(1)$.

\subsection{The Federated Audit Setting}

\begin{definition}[Federated Audit Problem]
\label{def:fed-problem}
$K$ parties each hold a private dataset $\cD_k = \{(x_i^{(k)}, y_i^{(k)})\}_{i=1}^{n_k}$,
$k = 1, \dots, K$.  A federated aggregator produces a global model $M_\Pi$
via protocol $\Pi$.  The auditor $\cA$ can access:
\begin{equation}\label{eq:fed-info}
    \cI_{\mathrm{Fed}} = \{S_k, \nabla S_k, H_{S,k}\}_{k=1}^K \cup \{M_\Pi\},
\end{equation}
where $S_k = \Cercis(\cD_k)$ is the Cercis Score computed locally by party
$k$, optionally encrypted before sharing.  The auditor \emph{cannot} access
the raw $(x, y)$ pairs.  The full-information audit statistic
$T_{\mathrm{full}}: (\cX \times \cY)^{n_{\mathrm{total}}} \to \R$ uses
$\cup_k \cD_k$.  The federated audit statistic
$T_{\mathrm{fed}}: \cI_{\mathrm{Fed}} \to \R$ uses only the audit
fingerprints.  The goal is to bound
$|\mathrm{Power}(T_{\mathrm{fed}}) - \mathrm{Power}(T_{\mathrm{full}})|$.
\end{definition}

\subsection{Situs Homomorphic Composition}

The central structural hypothesis enabling federated auditing is that the
Situs encoding admits a composition operator that preserves the geometry of
the union of distributions.

\begin{definition}[Situs Homomorphic Composition]
\label{def:situs-hom}
A binary operator $\oplus$ on Situs spaces is \textbf{homomorphic} with
respect to dataset union if there exists a function
$\phi: \R^K \to \R_{\geq 0}$ such that
\begin{equation}\label{eq:homomorphism}
    d_{\Situs}\!\big(\Situs(\cup_{k=1}^K \cD_k),\;
    \oplus_{k=1}^K \Situs(\cD_k)\big)
    \;\leq\; \phi(n_1, \dots, n_K),
\end{equation}
and $\phi(n_1, \dots, n_K) \to 0$ as $\min_k n_k \to \infty$.  That is, the
Situs encoding of the union is asymptotically approximated by the composition
of individual Situs encodings.
\end{definition}

\begin{remark}[Candidate Constructions]
Three candidate constructions for $\oplus$ are under investigation:
\begin{enumerate}[label=(\roman*)]
    \item \textbf{Wasserstein barycenter}: $\oplus_k \Situs_k =
          \argmin_{\nu} \sum_k \lambda_k W_2^2(\Situs_k, \nu)$ with
          entropic regularization~\citep{cuturi2013sinkhorn};
    \item \textbf{Gromov--Wasserstein fusion}: Align Situs spaces via
          Gromov--Wasserstein distance then average in the aligned space;
    \item \textbf{Kernel mean embedding}: $\oplus_k \Situs_k =
          \frac{1}{K}\sum_k \Phi(\Situs_k)$ in a reproducing kernel
          Situs space, where $\Phi$ is the kernel embedding.
\end{enumerate}
Each candidate has different trade-offs in computational cost, communication
overhead, and approximation rate $\phi$.  The existence of \emph{any}
$\oplus$ with $\phi = O(1/\sqrt{\min_k n_k})$ or better is an \textbf{open
problem}.  \openquest
\end{remark}

\subsection{Zero-Knowledge Audit Verification}

\begin{definition}[ZK Audit Proof]
\label{def:zk-proof}
A zero-knowledge proof system for Cercis Score submissions is a triple
$(\mathsf{Setup}, \mathcal{P}, \mathcal{V})$:
\begin{enumerate}[label=(\roman*)]
    \item \textbf{Completeness}: If party $k$ honestly computes
          $S_k = \Cercis(\cD_k)$ and commits via
          $c_k = \Commit(\cD_k; r_k)$, then
          $\mathcal{V}(c_k, S_k, \pi_k) = 1$ with probability
          $1 - \mathrm{negl}(\lambda)$, where
          $\pi_k \gets \mathcal{P}(\cD_k, r_k, S_k)$ and $\lambda$ is the
          security parameter.
    \item \textbf{Soundness}: For any $S'_k \neq \Cercis(\cD_k)$ and any
          polynomial-time adversary $\mathcal{P}^*$,
          $\Pr[\mathcal{V}(c_k, S'_k, \pi^*) = 1] \leq \mathrm{negl}(\lambda)$.
    \item \textbf{Zero-Knowledge}: There exists a simulator $\mathcal{S}$
          such that the view of any verifier interacting with $\mathcal{P}$
          is computationally indistinguishable from $\mathcal{S}(c_k, S_k)$.
\end{enumerate}
\end{definition}

The commitment $c_k$ binds party $k$ to its dataset $\cD_k$ without revealing
it.  The $\ZK$ proof $\pi_k$ attests that $S_k$ was correctly computed from
the committed $\cD_k$, without revealing $\cD_k$ itself.  This is the
cryptographic wedge that makes federated auditing possible: the auditor
verifies the \emph{computation} without seeing the \emph{data}.

% ============================================================
\section{Main Results}
% ============================================================

\subsection{Federated Audit $\varepsilon$-Equivalence}

\begin{theorem}[Federated Audit Completeness]
\label{thm:fed-equiv}
Let $T_{\mathrm{full}}$ be the full-information audit statistic and
$T_{\mathrm{fed}}$ the federated audit statistic based on
$\cI_{\mathrm{Fed}}$.  Assume:
\begin{enumerate}[label=(\alph*)]
    \item \textbf{ZK Verification}: Each party $k$ submits
          $(c_k, S_k, \pi_k)$ satisfying Definition~\ref{def:zk-proof};
    \item \textbf{Situs Homomorphism}: There exists $\oplus$ satisfying
          Definition~\ref{def:situs-hom} with error
          $\phi = \phi(n_1, \dots, n_K)$.
\end{enumerate}
Let $\delta \geq 0$ be the total differential privacy budget across all
parties (with $\delta = 0$ if no DP mechanism is employed).  Then
\begin{equation}\label{eq:fed-equiv-bound}
    \big|\mathrm{Power}(T_{\mathrm{fed}}) - \mathrm{Power}(T_{\mathrm{full}})\big|
    \;\leq\; \varepsilon(K, \delta, \phi),
\end{equation}
where
\begin{equation}\label{eq:epsilon-def}
    \varepsilon(K, \delta, \phi) =
    C_1 \cdot \phi
    + C_2 \cdot \sqrt{\frac{K \cdot \delta}{n_{\mathrm{total}}}}
    + C_3 \cdot \frac{1}{\sqrt{\min_k n_k}}
    + C_4 \cdot \mathrm{negl}(\lambda),
\end{equation}
with $C_1, C_2, C_3, C_4$ problem-dependent constants,
$n_{\mathrm{total}} = \sum_k n_k$.  As $\min_k n_k \to \infty$,
$\delta \to 0$, and $\lambda \to \infty$, we have
$\varepsilon(K, \delta, \phi) \to 0$ provided $\phi \to 0$.
\end{theorem}

\begin{proof}[Proof Sketch]
The total error decomposes into four independent sources, each bounded
separately and combined via the triangle inequality on test power.

\textit{Source 1: Situs composition error.}
The full-information statistic $T_{\mathrm{full}}$ depends on the joint
empirical distribution of $\cup_k \cD_k$.  Under the Situs homomorphism
assumption~\eqref{eq:homomorphism}, the composed Situs encoding
$\oplus_k \Situs_k$ approximates $\Situs(\cup_k \cD_k)$ to within $\phi$.
Since the Cercis Score is $L$-Lipschitz with respect to the Situs
metric~\citep{scx2024cercis}, the induced error in $T_{\mathrm{fed}}$
is bounded by $L \cdot \phi$, yielding the $C_1 \cdot \phi$ term.

\textit{Source 2: DP distortion.}
If parties apply $(\varepsilon_{\mathrm{DP}}, \delta)$-differential privacy,
the privatization mechanism adds noise of variance
$\sigma^2 \propto K/\varepsilon_{\mathrm{DP}}^2$ by composition across $K$
parties.  The effect on test power is bounded via the Gaussian DP
hypothesis testing framework~\citep{dong2022gaussian}, yielding the
$C_2 \cdot \sqrt{K \delta / n_{\mathrm{total}}}$ term.  When no DP is
used ($\delta = 0$), this term vanishes identically.

\textit{Source 3: Finite-sample estimation.}
Each $S_k$ is estimated from $n_k$ samples.  By Hoeffding's inequality
applied to the Cercis Score (bounded in $[0, 1+\eta]$), the estimation
error is $O(1/\sqrt{n_k})$.  The worst-case across parties gives the
$C_3 \cdot 1/\sqrt{\min_k n_k}$ term.

\textit{Source 4: Cryptographic soundness error.}
The $\ZK$ proof system contributes a negligible soundness error
$\mathrm{negl}(\lambda)$, which can be driven exponentially small by
increasing the security parameter.

The asymptotic behavior follows directly: $\phi \to 0$ (Situs homomorphism),
$\delta \to 0$ (DP-free or vanishing privacy budget), $\min_k n_k \to \infty$
(large datasets), $\lambda \to \infty$ (strong cryptography).  Under these
limits, $\varepsilon \to 0$ and federated audit achieves full-data audit
power.  $\square$
\end{proof}

\begin{remark}[Status]
The $\varepsilon$-bound is \textbf{conditionally rigorous}: given the Situs
homomorphic composition operator $\oplus$, the derivation is a strict
consequence of the Lipschitz property of Cercis, the DP hypothesis testing
literature, and standard concentration inequalities.  The open component is
the \emph{construction} of $\oplus$, not the logical chain that follows
from it.  \conditionallyrigorous~/ \openquest
\end{remark}

\subsection{Information Lower Bound: The Cost of Unverified Noise}

\begin{theorem}[Information Lower Bound for Federated Audit]
\label{thm:info-lower}
Suppose there exists a party $k$ for which the auditor has \emph{no}
$\ZK$ verification of $S_k$'s correctness — i.e., party $k$ can submit
an arbitrary $\tilde{S}_k \neq \Cercis(\cD_k)$ without detection.  Let
$\eta_k$ be the unknown label noise parameter of $\cD_k$.  Then for any
federated audit statistic $T_{\mathrm{fed}}$,
\begin{equation}\label{eq:lower-bound}
    \sup_{\cD}\;
    \big|\,\E[T_{\mathrm{fed}}] - \E[T_{\mathrm{full}}]\,\big|
    \;\geq\; c \cdot \eta_k \cdot \frac{n_k}{n_{\mathrm{total}}},
\end{equation}
where $c > 0$ is a universal constant and the supremum is taken over all
dataset configurations $\cD = (\cD_1, \dots, \cD_K)$ with fixed marginal
sizes.
\end{theorem}

\begin{proof}
The proof chains Theorem~3 with the data processing inequality.

\textit{Step 1: Reduction to Theorem~3.}
Without $\ZK$ verification, party $k$ controls its submitted $S_k$.  Consider
a strategy where party $k$ replaces its true dataset $\cD_k$ with a
\emph{confounding distribution} $\tilde{\cD}_k$ that: (a)~matches the
marginal statistics visible to the auditor (i.e., produces a Cercis Score
distribution consistent with some valid dataset), but (b)~inflates the
effective noise parameter from $\eta_k$ to $\tilde{\eta}_k$.  By Theorem~3,
the auditor cannot distinguish between a genuinely high-noise dataset and
one that has been adversarially manipulated, because both produce
observationally equivalent Cercis Score distributions when $\eta$ is unknown.

\textit{Step 2: Information-theoretic gap.}
Let $\Delta = \E[T_{\mathrm{fed}}] - \E[T_{\mathrm{full}}]$.  By the data
processing inequality,
\begin{equation}\label{eq:dpi}
    \I(T_{\mathrm{full}}; \cD_k) \;\geq\; \I(T_{\mathrm{fed}}; \cD_k),
\end{equation}
since $T_{\mathrm{fed}}$ is a (possibly randomized) function of the
manipulated submission.  The mutual information gap translates to an
estimation gap via the I-MMSE relationship~\citep{guo2005mutual}:
\begin{equation}\label{eq:immse}
    \E[(T_{\mathrm{full}} - \E[T_{\mathrm{full}} \mid \cD_k])^2]
    \leq \E[(T_{\mathrm{fed}} - \E[T_{\mathrm{fed}} \mid \cD_k])^2]
    + \frac{1}{2}\,\I(T_{\mathrm{full}}; \cD_k \mid T_{\mathrm{fed}}).
\end{equation}

\textit{Step 3: Worst-case construction.}
Construct the worst-case dataset $\cD^*$ by concentrating all label noise in
$\cD_k$ ($\eta_k \to 1$) while keeping $\eta_j = 0$ for $j \neq k$.  Under
this configuration, the $n_k$ samples whose quality is unverifiable contribute
a bias proportional to their share of the total data.  The adversary can tune
$\tilde{S}_k$ to maximize $|\E[T_{\mathrm{fed}}] - \E[T_{\mathrm{full}}]|$,
achieving the lower bound $c \cdot n_k / n_{\mathrm{total}}$.  Scaling by
the actual noise level $\eta_k$ yields~\eqref{eq:lower-bound}.  $\square$
\end{proof}

\begin{corollary}[ZK Necessity]
\label{cor:zk-necessity}
To achieve $\varepsilon$-equivalence with $\varepsilon \to 0$ as
$n_{\mathrm{total}} \to \infty$, $\ZK$ verification (or an equivalent
commitment-verification mechanism) is \emph{necessary} — not merely
sufficient.  Without it, the lower bound~\eqref{eq:lower-bound} persists
as a constant fraction regardless of sample size.
\end{corollary}

\begin{proof}
The bound~\eqref{eq:lower-bound} depends only on the \emph{proportion}
$n_k / n_{\mathrm{total}}$, not on the absolute scale.  No amount of data
from verified parties can compensate for an unverified party: the bias is
additive and non-vanishing.  $\square$
\end{proof}

\begin{remark}
Theorem~\ref{thm:info-lower} is \textbf{rigorous} — it follows directly
from Theorem~3 and the data processing inequality, requiring no additional
assumptions.  It provides the theoretical imperative for cryptographic
verification in federated SCX deployments.  \rigorous
\end{remark}

\subsection{ZK Protocol Construction Sketch}

While a full cryptographic implementation is beyond this paper's scope, we
sketch the architecture of a practical $\ZK$ proof for Cercis Score
computation.

\begin{algorithm}[tb]
\caption{ZK-Verified Cercis Score Submission (Party $k$)}
\label{alg:zk-cercis}
\begin{algorithmic}[1]
\Require Dataset $\cD_k$, security parameter $\lambda$, Situs spec
\Ensure $(c_k, S_k, \pi_k)$ — commitment, score, ZK proof
\State $r_k \gets \{0,1\}^\lambda$ \Comment{Sample randomness}
\State $c_k \gets \Commit(\cD_k; r_k)$ \Comment{Pedersen or Merkle commitment}
\State $S_k \gets \Cercis(\cD_k)$ \Comment{Compute Cercis Score locally}
\State $\mathcal{C} \gets \mathrm{Circuit}(\Cercis, \Commit)$
      \Comment{Arithmetic circuit encoding $\Cercis$ computation}
\State $\pi_k \gets \mathrm{Prove}(\mathcal{C}, (c_k, r_k, \cD_k), S_k)$
      \Comment{ZK-SNARK or Bulletproofs}
\State \Return $(c_k, S_k, \pi_k)$
\end{algorithmic}
\end{algorithm}

The arithmetic circuit $\mathcal{C}$ encodes three constraints:
(i)~$c_k$ is a valid commitment to $\cD_k$ with opening $r_k$;
(ii)~$S_k$ equals the Cercis Score computed from the committed $\cD_k$;
(iii)~the Cercis computation uses the canonical operator as specified by
the SCX framework.  The circuit size scales as
$O(n_k \cdot \dim(\cX) \cdot \log(1/\varepsilon_{\mathrm{prec}}))$, where
$\varepsilon_{\mathrm{prec}}$ is the arithmetic precision.

For large datasets, several mitigations apply: (i)~incremental Cercis
computation with recursive $\ZK$ proofs — each batch update produces a small
proof, and a final recursive proof composes them; (ii)~approximate Cercis
with bounded error, trading a small increase in $\varepsilon$ for reduced
circuit depth; (iii)~hardware-assisted Trusted Execution Environments (TEEs)
as a pragmatic alternative where TEE security guarantees suffice.

% ============================================================
\section{Telegraph Public Good Extension}
\label{sec:telegraph}

If the ``Telegraph Public Good'' — a decentralized data infrastructure
layer — extends to the federated setting, the architecture simplifies
substantially.  Instead of each party generating a $\ZK$ proof at audit
time, the storage layer itself provides a commitment to each $\cD_k$'s
contents, indexed by a Situs hash:

\begin{definition}[Situs Hash]
The Situs hash of a dataset $\cD_k$ is
$\mathcal{H}_{\Situs}(\cD_k) = \mathsf{CRH}(\Situs(\cD_k))$, where
$\mathsf{CRH}$ is a collision-resistant hash function acting on a
canonical serialization of the Situs encoding.
\end{definition}

Under this extension:
\begin{enumerate}[label=(\roman*)]
    \item Each party commits $\cD_k$ to the storage layer, which publishes
          $\mathcal{H}_{\Situs}(\cD_k)$;
    \item The auditor requests a \textbf{Proof of Storage} (PoS) that the
          data corresponding to $\mathcal{H}_{\Situs}(\cD_k)$ exists and is
          retrievable;
    \item The Cercis Score $S_k$ is computed \emph{by the auditor} (or within
          a Trusted Execution Environment) directly from the storage layer,
          eliminating the need for per-party $\ZK$ proofs of computation.
\end{enumerate}

This tightens the bound~\eqref{eq:fed-equiv-bound} by removing the
cryptographic soundness error $C_4 \cdot \mathrm{negl}(\lambda)$ and
reducing the finite-sample error (the auditor computes $S_k$ from the full
committed data, not from a party's summary).  The remaining error is purely
statistical: $C_3 \cdot 1/\sqrt{\min_k n_k}$.  Whether this architecture is
practically realizable depends on the design of the decentralized storage
layer — an \textbf{open problem} with a clear and well-defined theoretical
path.  \openquest

% ============================================================
\section{Discussion}
% ============================================================

\subsection{The Audit Paradox Resolved}

The federated audit paradox — auditing what you cannot see — is resolved not
at the cryptographic layer, but at the \textbf{information architecture}
layer.  The Situs encoding functions as an \emph{audit fingerprint}: it
preserves the statistical structure necessary for audit conclusions (noise
rates, quality rankings, gating decisions) while discarding the
reconstructability of raw data.  This is the fundamental wedge that makes
federated auditing possible.

The $\ZK$ proof (or its Telegraph equivalent) is the mechanism by which the
auditor verifies that the fingerprint is genuine.  Without this verification
layer, Theorem~\ref{thm:info-lower} demonstrates that the fingerprint can
be forged, and the audit conclusion is only as reliable as the least
verifiable party.

\subsection{Implications for Federated Learning Infrastructure}

Our results suggest a design principle for privacy-preserving ML
infrastructure: \textbf{additive auditability}.  Secure Aggregation, DP, and
MPC each add a layer of privacy protection; none adds a layer of
auditability.  The Situs homomorphic composition operator, if constructible,
would provide the missing layer — a mathematical guarantee that the union of
privately-held datasets can be audited \emph{as if} the auditor had seen
all the data, while provably preventing the auditor from \emph{actually}
seeing it.

This shifts the federated learning conversation from ``how do we prevent
leakage'' to ``how do we enable verification.''  Both are necessary; only
the former is currently solved.

\subsection{The Situs Homomorphism as a Structural Bottleneck}

The central obstacle — the construction of $\oplus$ — is not merely a
technical gap.  It probes a foundational question: is the Situs encoding
\emph{compositional}, i.e., does the Situs of a union decompose into a
function of the Situs encodings of the parts?  If the answer is negative,
then federated auditing in SCX requires a fundamentally different approach
— perhaps interactive protocols where the auditor sequentially queries
parties, or hierarchical Situs encodings that trade compositionality for
approximation fidelity.

% ============================================================
\section{Conclusion}
% ============================================================

We have formalized the Federated Audit Problem and established the conditions
under which SCX auditing extends to multi-party settings where raw data is
private.  Three results define the landscape:

\begin{enumerate}[label=(\arabic*)]
    \item The $\varepsilon$-equivalence theorem establishes that, under Situs
          homomorphic composition and $\ZK$ verification, federated audit
          power approximates full-data audit power with error that vanishes
          in the large-sample, zero-DP limit.  \conditionallyrigorous~/
          \openquest~(Situs homomorphism)
    \item The information lower bound proves that $\ZK$ verification (or
          equivalent commitment verification) is \emph{necessary} — without
          it, the audit power loss is bounded below by a constant fraction
          of the unverified noise.  \rigorous
    \item The Telegraph Public Good extension points toward a simpler
          architecture where the storage layer itself provides the
          commitment, eliminating per-party $\ZK$ proofs.  \openquest
\end{enumerate}

The resolution of the Situs homomorphic composition problem is the critical
next step.  Its solution — or the proof of its impossibility — will determine
whether the SCX Audit Sword can extend intact into the federated world, or
whether federated auditing requires a fundamentally different theoretical
apparatus.

% ============================================================
\bibliographystyle{plainnat}
\begin{thebibliography}{30}

\bibitem{scx2024cercis}
SCX Working Group.
\newblock ``The SCX Axiomatic Framework: Cercis, Situs, and Tiresias
Operators,''
\newblock \emph{Technical Report}, 2024.

\bibitem{bonawitz2017practical}
K.~Bonawitz, V.~Ivanov, B.~Kreuter, A.~Marcedone, H.~B.~McMahan, S.~Patel,
D.~Ramage, A.~Segal, and K.~Seth.
\newblock ``Practical secure aggregation for privacy-preserving machine
learning,''
\newblock in \emph{ACM CCS}, 2017.

\bibitem{abadi2016deep}
M.~Abadi, A.~Chu, I.~Goodfellow, H.~B.~McMahan, I.~Mironov, K.~Talwar, and
L.~Zhang.
\newblock ``Deep learning with differential privacy,''
\newblock in \emph{ACM CCS}, 2016.

\bibitem{evans2018pragmatic}
D.~Evans, V.~Kolesnikov, and M.~Rosulek.
\newblock ``A pragmatic introduction to secure multi-party computation,''
\newblock \emph{Foundations and Trends in Privacy and Security},
2(2--3):70--246, 2018.

\bibitem{goldwasser1989knowledge}
S.~Goldwasser, S.~Micali, and C.~Rackoff.
\newblock ``The knowledge complexity of interactive proof systems,''
\newblock \emph{SIAM Journal on Computing}, 18(1):186--208, 1989.

\bibitem{groth2016size}
J.~Groth.
\newblock ``On the size of pairing-based non-interactive arguments,''
\newblock in \emph{EUROCRYPT}, 2016.

\bibitem{bunz2018bulletproofs}
B.~B\"unz, J.~Bootle, D.~Boneh, A.~Poelstra, P.~Wuille, and G.~Maxwell.
\newblock ``Bulletproofs: Short proofs for confidential transactions and
more,''
\newblock in \emph{IEEE S\&P}, 2018.

\bibitem{dong2022gaussian}
J.~Dong, A.~Roth, and W.~J.~Su.
\newblock ``Gaussian differential privacy,''
\newblock \emph{Journal of the Royal Statistical Society: Series B},
84(1):3--37, 2022.

\bibitem{guo2005mutual}
D.~Guo, S.~Shamai, and S.~Verd\'u.
\newblock ``Mutual information and minimum mean-square error in Gaussian
channels,''
\newblock \emph{IEEE Transactions on Information Theory}, 51(4):1261--1282,
2005.

\bibitem{cuturi2013sinkhorn}
M.~Cuturi.
\newblock ``Sinkhorn distances: Lightspeed computation of optimal transport,''
\newblock in \emph{NeurIPS}, 2013.

\bibitem{kairouz2021advances}
P.~Kairouz, H.~B.~McMahan, et al.
\newblock ``Advances and open problems in federated learning,''
\newblock \emph{Foundations and Trends in Machine Learning},
14(1--2):1--210, 2021.

\bibitem{cover1999elements}
T.~M.~Cover and J.~A.~Thomas.
\newblock \emph{Elements of Information Theory}, 2nd ed.
\newblock Wiley, 2006.

\bibitem{villani2008optimal}
C.~Villani.
\newblock \emph{Optimal Transport: Old and New}.
\newblock Springer, 2008.

\end{thebibliography}

\end{document}
